% Lösungen Einheit 1 von 2
\loesungskopf{le1}{Einheit 1 von 2 · Streifen- und Kreisdiagramm}
\erg{12}{a) 28 \quad b) 250 \quad c) 180 \quad d) 640 \quad e) 24 (alle, die abgestimmt haben, auch die ungültigen Zettel)}
\erg{13}{a) Winkel \quad b) Anteil \quad c) Winkel \quad d) Anteil \quad e) Anteil \quad f) Winkel}
\erg{14}{a) 4\,cm \quad b) 7\,cm \quad c) 2\,cm \quad d) 9\,cm}
\erg{15}{a) $6{,}5$\,cm (65\,mm) \quad b) Rad $4{,}5$\,cm, Bus 4\,cm, zu Fuß $1{,}5$\,cm \quad c) kein Frühstück $30\,\%$; Müsli 5\,cm, Brot 2\,cm, kein Frühstück 3\,cm \quad d) Brötchen 4\,cm, Getränke 3\,cm, Müsliriegel $2{,}5$\,cm, Kekse $0{,}5$\,cm (5\,mm); zusammen 10\,cm, jeder Abschnitt beschriftet}
\erg{16}{a) $180^\circ$ \quad b) $90^\circ$ \quad c) $36^\circ$ \quad d) $270^\circ$ \quad e) $54^\circ$ \quad f) $0{,}05 \cdot 360^\circ = 18^\circ$ \quad g) $45^\circ$ \quad h) $\frac{3}{8} \cdot 360^\circ = 135^\circ$ \quad i) $13 : 40 = 0{,}325$; $117^\circ$ \quad j) $4 : 13 \cdot 360^\circ = 110{,}8\ldots^\circ \approx 111^\circ$}
\erg{17}{a) $144^\circ$ \quad b) Tee $180^\circ$, Kakao $162^\circ$, Milch $18^\circ$}
\erg{18}{a) Sonstiges $100\,\% - 42{,}5\,\% - 27{,}5\,\% - 18{,}5\,\% = 11{,}5\,\%$; Sport $153^\circ$, Musik $99^\circ$, Lesen $66{,}6^\circ \approx 67^\circ$, Sonstiges $41{,}4^\circ \approx 41^\circ$; Summe $360^\circ$ \quad b) vier Sektoren mit diesen Winkeln, jeder beschriftet}
\erg{19}{a) A: die Hälfte, $50\,\%$ \quad b) B: ein Viertel, $25\,\%$ \quad c) D: ein Drittel, $\approx 33\,\%$ \quad d) F: ein Sechstel, $\approx 17\,\%$}
\erg{20}{a) im Uhrzeigersinn ab oben: Handy $25\,\%$, Bücher $55\,\%$, Fernsehen $20\,\%$ \quad b) im Uhrzeigersinn ab oben: Unterhaltung $34\,\%$, Waschen $10\,\%$, Kühlen $24\,\%$, Sonstiges $14\,\%$, Kochen $18\,\%$ (der größte Sektor zum größten Anteil) \quad c) $100 - 34 - 24 - 18 - 10 = 14\,\%$ \quad d) $72 : 3200 \cdot 360^\circ = 8{,}1^\circ$}
\erg{21}{a) Bei „zu Fuß“ ist der Prozentwert als Grad eingetragen; richtig $0{,}15 \cdot 360^\circ = 54^\circ$ \quad b) Kakao $198^\circ$, Tee $72^\circ$, Wasser $90^\circ$ \quad c) Timo nimmt die Anzahl 12 als Prozentsatz; richtig $12 : 300 = 0{,}04 = 4\,\%$, $0{,}04 \cdot 360^\circ = 14{,}4^\circ$ \quad d) $9 : 450 = 0{,}02$; $7{,}2^\circ$}
\erg{22}{a) $100\,\%$ sind der ganze Kreis mit $360^\circ$, also ist $1\,\%$ der hundertste Teil: $360^\circ : 100 = 3{,}6^\circ$. \quad b) Nein: Der Winkel hängt nur vom Anteil ab. $50\,\%$ sind in beiden Klassen die Hälfte des Kreises, $180^\circ$.}
\erg{23}{a) Fußball $45\,\%$, Schwimmen $30\,\%$, Radfahren $40\,\%$, Tanzen $20\,\%$ \quad b) $135\,\%$ \quad c) Nein: Die Anteile ergeben mehr als $100\,\%$, weil Kinder mehrere Sportarten nennen durften; ein Kreis kann nur $100\,\%$ aufteilen. Passend ist ein Säulen- oder Balkendiagramm.}

\noindent Lösungsgrafiken 17\,b) und 18\,b):\par
\noindent\kreisdiagramm[1.2]{Tee/50, Kakao/45, Milch/5}\hspace{1.5cm}\kreisdiagramm[1.2]{Sport/42.5, Musik/27.5, Lesen/18.5, Sonstiges/11.5}
